% ECONOMICS_PROBLEM: {"id": "C73-145", "title": "Learning with persistent states", "jel_code": "C73", "source_id": "OP-0285"}
% !TeX program = xelatex
\documentclass[11pt,a4paper]{article}
\usepackage[margin=23mm]{geometry}
\usepackage{fontspec}
\setmainfont{Latin Modern Roman}
\usepackage{amsmath,amssymb,mathtools}
\usepackage{xcolor,enumitem,booktabs,longtable,array,xurl,hyperref}
\definecolor{muted}{HTML}{53636C}
\hypersetup{colorlinks=true,urlcolor=blue,linkcolor=blue}
\setlength{\parindent}{0pt}
\setlength{\parskip}{5pt}
\newcommand{\R}{\mathbb R}
\newcommand{\E}{\mathbb E}
\newcommand{\N}{\mathbb N}
\newcommand{\ind}{\mathbf 1}
\newcommand{\Var}{\operatorname{Var}}
\newcommand{\Cov}{\operatorname{Cov}}
\newcommand{\supp}{\operatorname{supp}}
\DeclareMathOperator*{\argmin}{arg\,min}
\DeclareMathOperator*{\argmax}{arg\,max}
\newcommand{\entrymeta}[3]{{\sffamily\small\color{muted}\textbf{\texttt{#1}}\quad|\quad #2\par #3\par}}
\newcommand{\Problemref}[1]{\texttt{#1}}
\newcommand{\JELref}[2]{\texttt{#2} (JEL \texttt{#1})}
\begin{document}
\section*{Learning with persistent states}
\textbf{Problem ID:} \texttt{C73-145}\quad \textbf{JEL:} \texttt{C73}\par
% BEGIN ECONOMICS STATEMENT
\label{problem:OP-0285}
{\sffamily\small\color{muted}Original subject group: Game theory and strategic interaction\par}
\label{definitions:problem:30007567}
\textbf{Audit classification:} Published research agenda. Tardos’s 2026 ICM lecture explicitly identifies carry-over learning as a research area. The uniform welfare inequality is editorial and is false for arbitrary learners.
\subsection*{Definitions and precise research target}
\textbf{Model and research target.} Consider a finite-state resource-allocation game with state space \(S\), finite action sets \(A_i\), reward \(r_i(s,a)\in[0,1]\), and transition kernel \(P(s'\mid s,a)\). States record budgets or other resources carried between periods. Fix initial law \(\nu\), a class \(\mathcal C\) of admissible kernels and reward systems, and a class \(\mathcal L\) of decentralized learning rules using explicitly specified feedback. Let \(W_T(L)=E_{\nu,L}[\sum_{t=1}^T\sum_i r_i(s_t,a_t)]\). Let \(W_T^*\) be the maximal expected welfare of a centralized controller with the same initial law, transitions, and resource constraints. Determine assumptions and behavioral performance criteria under which
\[
 W_T(L)\geq c\,W_T^*-o(T)
 \quad\text{for every game in }\mathcal C\text{ and every }L\in\mathcal L,
\]
with a nonzero constant \(c\) independent of the horizon. Define the little-order remainder uniformly over the stated class. Any regret comparator must simulate the state trajectory induced by an alternative policy; holding observed states fixed is generally inadequate when actions alter tomorrow's resources. Characterize which mixing, smoothness, information, or budget assumptions are necessary, and give tight failures outside the guarantee.

\emph{Definitions and maintained assumptions (editorial unless a source definition is named).} A policy or learner is a stochastic map from its available history to an action distribution; declare whether it observes state, own reward, opponent actions and transition outcomes. The centralized benchmark optimizes a joint policy using the same physical state dynamics and an explicitly stated information advantage. A class \(\mathcal L\) cannot include arbitrary always-wasteful rules if positive welfare guarantees are sought. Specify a behavioral requirement, for example policy regret against an admissible stationary policy class, with each comparator’s entire state trajectory simulated under that policy. Define a uniform remainder \(g(T)\ge0\) by \(g(T)/T\to0\), independent of the game and learner, and require \(W_T(L)\ge cW_T^*-g(T)\). To make this meaningful, state uniform bounds on state and action counts, rewards, mixing times and the policy-regret bound; unrestricted state-class variation can destroy uniformity. Finite-state existence or finite-horizon dynamic programming for one known instance is a baseline, not the claimed frontier.
\subsection*{Variants and assumption changes}
\begin{enumerate}
\item \textbf{Uniform mixing (editorial).} Require every stationary joint policy to induce an irreducible chain with a common mixing-time bound. Seek welfare guarantees under a specifically defined policy-regret condition and known reward smoothness.
\item \textbf{Resource traps (editorial).} Permit an absorbing depleted-resource state. Construct tight examples showing whether exploration or irreversible budget use prevents positive uniform guarantees, then test restoration by replenishment or reset assumptions.
\end{enumerate}
\subsection*{References and source audit}
\begin{enumerate}
\item Learning outcomes in dynamic games; repeated auctions with budgets. Supports the stated research area; exact displayed specification is editorial. \url{https://epubs.siam.org/doi/full/10.1137/25M1808064}
\end{enumerate}
\begin{enumerate}\raggedright
\item\hypertarget{definitions:source:30007567:1}{} Éva Tardos. 2026. \emph{Learning and the Price of Anarchy in Games}. Learning outcomes in dynamic games; Packet routing as a dynamic game; Theorem8.
\url{https://epubs.siam.org/doi/full/10.1137/25M1808064}.
DOI: \url{https://doi.org/10.1137/25M1808064}.
\end{enumerate}

\subsection*{Common conventions: Source mathematical conventions}

These conventions apply to each entry unless explicitly replaced.

\textbf{Models and identification.} A primitive model \(m\) specifies all probability laws, feasible choices, information, equilibrium and measurement rules used by its target. The admissible class \(\mathcal M\) is fixed independently of outcomes. If \(P_O(m)\) is its observable probability law and \(\tau(m)\) the target, its identified set at observable law \(P\) is
\[
 \mathcal I_\tau(P)=\{\tau(m):m\in\mathcal M,\ P_O(m)=P\}.
\]
Point identification means that this set is a singleton. An empty set signals incompatibility of data and assumptions. Coordinate bounds need not describe the joint set. Bounds are sharp only if every retained target is attainable or arbitrarily approachable by compatible models. Finite bounds require explicit restrictions on unobserved quantities; unrestricted confounding, hidden wealth, extrapolation or arbitrary equilibrium selection can yield uninformative sets.

\textbf{Causal interventions.} A potential outcome \(Y_i(p)\) is the outcome under a complete policy regime \(p\), including timing, financing, information and market spillovers. The target population and units are fixed before treatment. With interference, use the complete assignment vector or a justified exposure mapping; individual no-interference notation alone cannot identify an economy-wide intervention. A difference of mean potential outcomes does not require the cross-world joint distribution. Conditioning on post-treatment migration, survival, enrollment or employment changes the estimand and needs separate justification.

\textbf{Statistical statements.} An estimator, confidence set, forecast or test specifies a sampling law, model class and loss/coverage criterion. Uniform \((1-\alpha)\)-coverage means \(\inf_{m\in\mathcal M}\Pr_m\{\tau(m)\in C_n\}\geq1-\alpha\), or an explicitly asymptotic version. The whole parameter space trivially has coverage, so informativeness requires a stated width, risk or power objective. A forecasting improvement is relative to a named benchmark, information set and evaluation population.

\textbf{Equilibrium and optimization.} An equilibrium comprises feasible plans, optimality, beliefs where needed, budget constraints and all market/resource clearing. The primitives or a stated benchmark define each correspondence \(\mathcal E(p;m)\). Nonemptiness is assumed or proved, never inferred just from compact policy sets. A selected-equilibrium target uses a declared selection rule; a robust target quantifies over all permitted equilibria. Use suprema or infima unless attainment is established. An \(\varepsilon\)-optimal policy is feasible and within \(\varepsilon\) of the finite optimum value. Report an empty feasible set separately.

\textbf{Welfare and accounting.} Normative utility, interpersonal normalization, social weights, numeraire, discounting and the use of public receipts are primitives. Behavioral utility need not coincide with the welfare criterion. Household welfare includes ownership income and rents once; adding profits again to compensating variation generally double-counts them. A monetary equivalent is defined by an explicit transfer and price convention, with monotonicity and range conditions. Average effects, mechanism contributions, transfers and welfare losses are different targets. Infinite sums require convergence and dynamic feasible plans require terminal or no-Ponzi conditions.

\textbf{Source versions.} A working-paper date, revision date and journal publication year are distinguished. A DOI for a journal version is not automatically the DOI of an earlier manuscript. Locators refer to the stated version and printed pages where specified. A metadata-only check does not verify a claim about a full-text conclusion. Every entry's source subsection narrows the provenance claim accordingly.
% END ECONOMICS STATEMENT
\end{document}
